\documentclass[11pt]{article}

\usepackage[margin=1in]{geometry}
\usepackage{amsmath}
\usepackage{booktabs}
\usepackage[hidelinks]{hyperref}

\title{Separation Pretraining for Convolutional Networks}
\author{AI in Science Lab}
\date{Work in Progress}

\begin{document}

\maketitle

\begin{abstract}
This experiment studies whether a convolutional neural network can learn more
effectively when its filters are first encouraged to produce distinct
responses. Before standard supervised training, selected convolutional layers
are optimized with a separation loss based on cosine similarity. We compare a
baseline with four variants that apply separation pretraining to progressively
more layers. The implementation and evaluation pipeline are complete; the
full comparison across datasets and random seeds remains to be run.
\end{abstract}

\section{Motivation}

Different convolutional filters are useful when they detect different
features.  However, filters may initially learn similar responses and therefore
use model capacity inefficiently.  Our working hypothesis is that briefly
encouraging filters to behave differently before classification training may
give the network a better starting point and help it learn faster.

\section{Method}

For each image in a batch and each selected convolutional layer, we sample
pairs of spatial patches from the feature map entering that layer. Each patch
has the same height and width as the layer's convolutional kernel. Applying all
filters in the layer to a patch produces a response vector $r$ with one entry
per output channel. For two sampled patches with response vectors $r_i$ and
$r_j$, their similarity is

\begin{equation}
    c_{ij} = \frac{r_i^\top r_j}
    {\lVert r_i \rVert_2\lVert r_j \rVert_2}.
\end{equation}

The loss for the pair is the squared cosine similarity,

\begin{equation}
    \mathcal{L}_{\mathrm{sep}}(i,j) = c_{ij}^{2}.
\end{equation}

Squaring penalizes both positive and negative alignment.  The minimum occurs
when the two response vectors are orthogonal, so the objective encourages
different responses without requiring one filter to become the negative of
another. The final separation loss is the mean over sampled patch pairs, input
images, and all selected layers.

During this phase, inputs to later convolutional layers are detached from the
gradient graph.  Consequently, a layer's separation signal updates that layer
rather than propagating through all earlier layers.  After separation
pretraining, the network is trained normally with cross-entropy loss.  The two
losses are used in separate phases and are not mixed.

\section{Model and Data}

The model is a compact convolutional classifier with four convolutional layers.
Their default widths are 32, 64, 128, and 256 channels. Each convolution uses a
$3\times3$ kernel and a rectified linear activation. Max pooling is applied
after the second, third, and fourth layers, followed by a linear classifier.

The experiment supports MNIST and CIFAR-10.  Both contain ten classes, while
their input shapes differ: MNIST uses $1\times28\times28$ grayscale images and
CIFAR-10 uses $3\times32\times32$ color images.  Dataset-specific normalization
is applied before training.

\section{Experimental Design}

We compare five variants that differ only in which convolutional layers receive
separation pretraining.

\begin{table}[h]
    \centering
    \begin{tabular}{ll}
        \toprule
        Variant & Pretrained convolutional layers \\
        \midrule
        Baseline & None \\
        Layer 1 & First layer \\
        Layers 1--2 & First and second layers \\
        Layers 1--3 & First, second, and third layers \\
        Layers 1--4 & All four layers \\
        \bottomrule
    \end{tabular}
    \caption{Separation-pretraining variants.}
\end{table}

The current default configuration uses AdamW, a batch size of 128, and a
learning rate of $10^{-3}$ for both phases. Separation pretraining lasts three
epochs and samples 16 patch pairs per image and selected layer by default.
Standard classification training then runs for 20 epochs. The planned grid
evaluates CIFAR-10 using the baseline, two-layer, and four-layer variants at
seed 0. It varies batch size, the learning rates for both training phases,
separation duration (one or five epochs), and patch-pair count (8 or 32),
producing 96 runs.

Training and validation loss, classification accuracy, learning rates, and the
average absolute cosine similarity are recorded with Weights \& Biases.  The
main comparison metric is validation accuracy.  Learning curves will also be
examined to test the more specific claim that separation pretraining improves
the speed of learning.

\section{Current Status and Next Steps}

The dataset loaders, model, separation objective, two-stage training procedure,
configuration, checkpoint support, and experiment logging are implemented.
No aggregate numerical results are included yet because the complete grid has
not been run and no result artifacts are currently stored with the experiment.

The next step is to run the full dataset--variant--seed grid.  Results should be
summarized using the mean and standard deviation across seeds.  In addition to
final validation accuracy, comparing validation accuracy over training time
will show whether the proposed initialization provides an early optimization
benefit.  Test-set evaluation should be reserved for the final selected setup.

\end{document}
